{{دیوان
| عنوان = ابی سِینیا
| ماخذ = ویکی ماخذ
| ماخذ_لنک = https://ur.wikisource.org/wiki/%D8%A7%D8%A8%DB%8C_%D8%B3%D9%90%DB%8C%D9%86%DB%8C%D8%A7
| مقابلہ = اقبال سائبر لائبریری (اقبال اکادمی پاکستان)
| مقابلہ_لنک = https://www.iqbalcyberlibrary.net
}}
=== (۱۸ اگست ۱۹۳۵ء) ===

<poem>
یورپ کے کرگسوں کو نہیں ہے ابھی خبر
ہے کتنی زہر ناک ابی سِینیا کی لاش

ہونے کو ہے یہ مُردۂ دیرینہ قاش قاش!

تہذیب کا کمال شرافت کا ہے زوال
غارت گری جہاں میں ہے اقوام کی معاش

ہر گُرگ کو ہے برّۂ معصوم کی تلاش!

اے وائے آبُروئے کلیسا کا آئِنہ
روما نے کر دیا سرِ بازار پاش پاش

پیرِ کلیسیا! یہ حقیقت ہے دلخراش!
</poem>
